\subsection{Gromov--Hausdorff convergence and compactness}\label{subsec:preliminaries-metric}
\paraid{p-9622f56720d9}
We begin with the correspondence formula for the distance and compactness criteria.

\paraid{p-f7c8be172ef5}
For a metric space
$\MetricSpace{\BaseCarrier}{\MetricSymbol}$,
a point
$\BasePoint\in\BaseCarrier$,
and a radius
$\Radius>0$,
we write
\[
 \MetricBall(\BasePoint,\Radius;\MetricSymbol)
 =\{\ComparisonPoint\in\BaseCarrier\mid\MetricSymbol(\BasePoint,\ComparisonPoint)<\Radius\}.
\]
For a subset
$\BorelSubset\subset\BaseCarrier$,
we put
\[
 \dist_{\MetricSymbol}(\BasePoint,\BorelSubset)
 =\inf_{\ComparisonPoint\in\BorelSubset}\MetricSymbol(\BasePoint,\ComparisonPoint),
 \qquad \dist_{\MetricSymbol}(\BasePoint,\emptyset)=\infty.
\]
We denote by
$\CompactHyperspace(\BaseCarrier)$
the set of all nonempty compact subsets of
$\BaseCarrier$.
The \emph{Hausdorff distance} on this set is defined by
\[
 \HausdorffDistance{\MetricSymbol}\colon
 \CompactHyperspace(\BaseCarrier)\times\CompactHyperspace(\BaseCarrier)\to[0,\infty).
\]
For every
$\FirstCompactSet,\SecondCompactSet\in\CompactHyperspace(\BaseCarrier)$,
put
\[
 \HausdorffDistance{\MetricSymbol}(\FirstCompactSet,\SecondCompactSet)
 =\max\left\{
   \sup_{\BasePoint\in\FirstCompactSet}\dist_{\MetricSymbol}(\BasePoint,\SecondCompactSet),
   \sup_{\ComparisonPoint\in\SecondCompactSet}\dist_{\MetricSymbol}(\ComparisonPoint,\FirstCompactSet)
 \right\}.
\]
The \emph{Vietoris topology} on
$\CompactHyperspace(\BaseCarrier)$
has a basis consisting of the sets obtained as follows.
For an integer
$m\geq1$
and open subsets
$U_1,\ldots,U_m$
of
$\BaseCarrier$,
take
\[
 \left\{\FirstCompactSet\in\CompactHyperspace(\BaseCarrier)\;\middle|\;
 \FirstCompactSet\subset\bigcup_{i=1}^{m}U_i,
 \quad \FirstCompactSet\cap U_i\ne\emptyset\ (1\leq i\leq m)\right\}.
\]

\begin{theorem}[{\cite[Exercises~2.2.11(a),(b) and~3.2.9]{Beer1993}}]\label{thm:compact-hyperspace-topology}
\paraid{p-974030144cf7}
Let
$\MetricSpace{\BaseCarrier}{\MetricSymbol}$
be a metric space.
Then the Hausdorff distance induces the Vietoris topology on
$\CompactHyperspace(\BaseCarrier)$.
If
$\BaseCarrier$
is separable,
then
$\CompactHyperspace(\BaseCarrier)$
is separable.
\end{theorem}

\begin{lemma}[{\cite[Chapter~2, Proposition~5.3, p.~40]{Wicks1991}}]\label{lem:hausdorff-lipschitz-images}
\paraid{p-eb3ed5b8a76b}
Let
$\MetricSpace{\BaseCarrier}{\MetricSymbol}$
and
$\MetricSpace{\ComparisonCarrier}{\ComparisonMetric}$
be metric spaces,
let
$L\geq0$,
and let
$\Mapping\colon\BaseCarrier\to\ComparisonCarrier$
be
$L$-Lipschitz.
Then for all
$\FirstCompactSet,\SecondCompactSet\in\CompactHyperspace(\BaseCarrier)$,
we have
$\HausdorffDistance{\ComparisonMetric}
 (\Mapping(\FirstCompactSet),\Mapping(\SecondCompactSet))
 \leq L\cdot \HausdorffDistance{\MetricSymbol}(\FirstCompactSet,\SecondCompactSet)$.
\end{lemma}

% \paraid{p-109cc338209b}

\paraid{p-0a1ce4229ada}
For
$\ErrorTolerance>0$,
a subset
$\NetSet\subset\BaseCarrier$
is an
\emph{$\ErrorTolerance$-net}
if
\[
 \forall\BasePoint\in\BaseCarrier\quad
 \exists\NetPoint\in\NetSet\quad
 \MetricSymbol(\BasePoint,\NetPoint)<\ErrorTolerance.
\]
Every compact metric space admits a finite
$\ErrorTolerance$-net
for every
$\ErrorTolerance>0$.

\paraid{p-d7c66dfce89f}
A \emph{correspondence} between nonempty sets
$\BaseCarrier $
and
$\ComparisonCarrier $
is a subset of
$\BaseCarrier \times \ComparisonCarrier $
whose two coordinate projections are surjective.
We denote by
$\Cor(\BaseCarrier,\ComparisonCarrier)$
the set of all correspondences between
$\BaseCarrier$
and
$\ComparisonCarrier$.
For metric spaces
$\MetricSpace{\BaseCarrier }{\MetricSymbol }$
and
$\MetricSpace{\ComparisonCarrier }{\ComparisonMetric }$,
the distortion of a correspondence
$\Correspondence\in\Cor(\BaseCarrier,\ComparisonCarrier)$
is
\[
 \dis(\Correspondence )=
 \sup\left\{
   \abs{\MetricSymbol (\BasePoint _0,\BasePoint _1)-\ComparisonMetric (\ComparisonPoint _0,\ComparisonPoint _1)}
   \,\middle|\,
   (\BasePoint _0,\ComparisonPoint _0),(\BasePoint _1,\ComparisonPoint _1)\in \Correspondence
 \right\}.
\]
\begin{theorem}[{\cite[Theorem~7.3.25]{BuragoBuragoIvanov2001} and \cite[arXiv v1, Theorem~6.12]{Tuzhilin2020}}]\label{thm:gh-correspondence}
\paraid{p-fc061a808065}
For nonempty compact metric spaces
$\MetricSpace{\BaseCarrier }{\MetricSymbol }$
and
$\MetricSpace{\ComparisonCarrier }{\ComparisonMetric }$,
the Gromov--Hausdorff distance satisfies
\begin{equation}\label{eq:correspondence}
 \GHDistance(\MetricSpace{\BaseCarrier }{\MetricSymbol },\MetricSpace{\ComparisonCarrier }{\ComparisonMetric })=\frac12\inf_{\Correspondence\in\Cor(\BaseCarrier,\ComparisonCarrier)}\dis(\Correspondence).
\end{equation}
Equivalently,
it is the infimum of the Hausdorff distances between isometric copies
of
$\MetricSpace{\BaseCarrier }{\MetricSymbol }$
and
$\MetricSpace{\ComparisonCarrier }{\ComparisonMetric }$
in a common metric space.

\end{theorem}

\paraid{p-c3f214f2a9fe}
The next lemma realizes a prescribed correspondence in a common metric space
with a bound on the distances between corresponding points.
We use \Cref{lem:correspondence-embedding} in the proof of
\Cref{lem:common-gh-embedding}
to embed a convergent sequence of compact metric spaces into one compact metric space.

\begin{lemma}\label{lem:correspondence-embedding}
\paraid{p-77a82a051441}
Let
$\MetricSpace{\BaseCarrier}{\MetricSymbol}$
and
$\MetricSpace{\ComparisonCarrier}{\ComparisonMetric}$
be nonempty compact metric spaces,
let
$\Correspondence\subset\BaseCarrier\times\ComparisonCarrier$
be a correspondence,
and let
$\SmallError>0$
satisfy
$\dis(\Correspondence)\leq2\SmallError$.
Then there is a metric
$\AmbientMetric$
on the disjoint union
$\BaseCarrier\sqcup\ComparisonCarrier$
extending both metrics such that for every
$(\BasePoint,\ComparisonPoint)\in\Correspondence$
we have
\[
 \AmbientMetric(\BasePoint,\ComparisonPoint)\leq\SmallError.
\]
In particular,
$\HausdorffDistance{\AmbientMetric}(\BaseCarrier,\ComparisonCarrier)\leq\SmallError$.
\end{lemma}

\begin{proof}
\paraid{p-1dc4840c170b}
This follows from
\cite[proof of Theorem~7.3.25, pp.~257--258]{BuragoBuragoIvanov2001}.
\end{proof}

\begin{theorem}[{Consequence of \cite[arXiv v1, Theorem~7.18]{Tuzhilin2020}}]\label{thm:uniform-nets}
\paraid{p-b1b76db6a357}
If
$\MetricSpace{\BaseCarrier_\SequenceIndex}{\MetricSymbol_\SequenceIndex}
 \to\MetricSpace{\BaseCarrier}{\MetricSymbol}$
in
$\GHSpace$,
then for every
$\ErrorTolerance>0$
there is a positive integer
$\NetSize$
such that each
$\BaseCarrier_\SequenceIndex$
admits an
$\ErrorTolerance$-net
of cardinality at most
$\NetSize$.
\end{theorem}

\paraid{p-e041a46858f0}
Let
$\MetricSpace{\BaseCarrier}{\MetricSymbol}$
be a metric space, and let
$\SeparationScale>0$.
A subset
$S\subset\BaseCarrier$
is called
\emph{$\SeparationScale$-separated}
if
\[
 \forall x,y\in S,\qquad
 x\ne y\Longrightarrow\MetricSymbol(x,y)\ge\SeparationScale.
\]

\begin{corollary}\label{cor:uniform-separated-sets}
\paraid{p-f9bb5a038d1d}
Let
$\MetricSpace{\BaseCarrier_\SequenceIndex}{\MetricSymbol_\SequenceIndex}
 \to\MetricSpace{\BaseCarrier}{\MetricSymbol}$
in
$\GHSpace$,
and let
$\SeparationScale>0$.
Then there is a positive integer
$\NetSize$
such that every
$\SeparationScale$-separated
subset of every
$\BaseCarrier_\SequenceIndex$
has cardinality at most
$\NetSize$.
\end{corollary}

\begin{proof}
\paraid{p-0e6ff4e87f94}
By \Cref{thm:uniform-nets},
choose
$\SeparationScale/3$-nets
$\mathcal{N}_\SequenceIndex\subset\BaseCarrier_\SequenceIndex$
in all
$\BaseCarrier_\SequenceIndex$
with
$\Card(\mathcal{N}_\SequenceIndex)\leq\NetSize$.
For an arbitrary
$\SeparationScale$-separated subset
$S_\SequenceIndex\subset\BaseCarrier_\SequenceIndex$,
choose for each point of
$S_\SequenceIndex$
a point of
$\mathcal{N}_\SequenceIndex$
within distance
$\SeparationScale/3$.
This assignment is injective,
because two points with the same image would have distance less than
$2\SeparationScale/3$.
Consequently,
$\Card(S_\SequenceIndex)
 \leq \Card(\mathcal{N}_\SequenceIndex)\leq \NetSize$.
This proves the uniform cardinality bound.
\end{proof}

\paraid{p-47d23e54815b}
We use completeness and separability of the Gromov--Hausdorff space
$\GHSpace$.

\begin{theorem}[{\cite[arXiv v1, Corollary~7.24]{Tuzhilin2020}}]\label{lem:gh-basics}
\paraid{p-fa3b1c89e980}
The metric space
$\MetricSpace{\GHSpace}{\GHDistance}$
is complete and separable.
\end{theorem}

\paraid{p-2215668b1fa1}
Gromov constructs a common compact metric space admitting isometric embeddings
of every member of a family of compact metric spaces with uniformly bounded
diameters and uniformly bounded covering numbers at each positive radius
\cite[Section~6, pp.~64--65]{Gromov1981}.
See also \cite[arXiv v1, Theorem~7.19]{Tuzhilin2020}.
In the next theorem,
we also arrange Hausdorff convergence of the embedded spaces.

\begin{theorem}\label{lem:common-gh-embedding}
\paraid{p-fc625e3c11fb}
Every Cauchy sequence
$\{\MetricSpace{\BaseCarrier_\SequenceIndex}{\MetricSymbol_\SequenceIndex}\}_{\SequenceIndex\in\NonnegativeIntegers}$
in
$\GHSpace$
admits isometric embeddings into a compact metric space
$\MetricSpace{\AmbientCarrier}{\AmbientMetric}$
whose images converge in Hausdorff distance to a nonempty compact subset.
If its Gromov--Hausdorff limit is
$\MetricSpace{\BaseCarrier}{\MetricSymbol}$,
then we may also embed
$\MetricSpace{\BaseCarrier}{\MetricSymbol}$
into the same ambient space.
Identifying these spaces with their images, we obtain
\[
 \HausdorffDistance{\AmbientMetric}(\BaseCarrier_\SequenceIndex,\BaseCarrier)\longrightarrow0.
\]
\end{theorem}

\begin{proof}
\paraid{p-0e319197acba}
By \Cref{lem:gh-basics},
let
$\MetricSpace{\BaseCarrier}{\MetricSymbol}$
be the limit and put
\[
 \ErrorTolerance_\SequenceIndex
 =\GHDistance(\MetricSpace{\BaseCarrier_\SequenceIndex}{\MetricSymbol_\SequenceIndex},
              \MetricSpace{\BaseCarrier}{\MetricSymbol})+2^{-\SequenceIndex}.
\]
Using \Cref{thm:gh-correspondence,lem:correspondence-embedding}
and gluing along the common copy of
$\BaseCarrier$
\cite[author's manuscript, Theorem~27.2]{Villani2009},
we obtain a metric
$\AmbientMetric$
on the union of
$\BaseCarrier$
and all
$\BaseCarrier_\SequenceIndex$
extending their metrics and satisfying
\[
 \HausdorffDistance{\AmbientMetric}
 (\BaseCarrier_\SequenceIndex,\BaseCarrier)
 \leq\ErrorTolerance_\SequenceIndex\longrightarrow0.
\]

\paraid{p-82732ff2ca3f}
For every
$\SmallError>0$,
the tail lies in the
$\SmallError/2$-neighborhood of
$\BaseCarrier$.
A finite
$\SmallError/2$-net in
$\BaseCarrier$,
together with finite
$\SmallError$-nets in the finitely many remaining spaces,
forms a finite
$\SmallError$-net for the union.
The completion of this union is therefore compact and retains the displayed
Hausdorff bounds.
\end{proof}

\paraid{p-159a1141fda7}
We also recall the compactness criterion used for families of functions.

\begin{theorem}[{Arzel\`a--Ascoli, \cite[Theorem~6.26]{Muscat2024}}]\label{thm:arzela-ascoli}
\paraid{p-017c33e53ae1}
Let
$\MetricSpace{\BaseCarrier}{\MetricSymbol}$
be nonempty and compact,
and let
$\MetricSpace{\ComparisonCarrier}{\ComparisonMetric}$
be complete.
Equip the set
$\ContinuousFunctions(\BaseCarrier,\ComparisonCarrier)$
of continuous maps with the uniform metric
\[
 \UniformMetric(\Mapping,\SecondMapping)
 =\sup_{\BasePoint\in\BaseCarrier}
   \ComparisonMetric(\Mapping(\BasePoint),\SecondMapping(\BasePoint)).
\]
Then a family
$\FunctionFamily\subset\ContinuousFunctions(\BaseCarrier,\ComparisonCarrier)$
has compact closure if and only if its combined image
$\bigcup_{\Mapping\in\FunctionFamily}\Mapping(\BaseCarrier)$
is totally bounded and the family
$\FunctionFamily$
is equicontinuous.
Here equicontinuity means that for every
$\ErrorTolerance>0$
there is
$\SmallError>0$
such that for every
$\Mapping\in\FunctionFamily$
and every
$\BasePoint,\ComparisonPoint\in\BaseCarrier$
with
$\MetricSymbol(\BasePoint,\ComparisonPoint)<\SmallError$,
we have
\[
 \ComparisonMetric(\Mapping(\BasePoint),\Mapping(\ComparisonPoint))<\ErrorTolerance.
\]
\end{theorem}

\paraid{p-f1481eeb5719}
The cited theorem characterizes total boundedness in the uniform metric.
Completeness of the target
$\MetricSpace{\ComparisonCarrier}{\ComparisonMetric}$
implies completeness of the function space
$\ContinuousFunctions(\BaseCarrier,\ComparisonCarrier)$
with the uniform metric
$\UniformMetric$,
which yields the formulation in terms of compact closure above.
